%% Stand 18. Juli 2026
\subsection{Grundlegende Begriffe}
\label{KapitelTheorieGrundlagen}
Das Fundament der nachfolgenden Überlegungen bildet die 
Annahme, dass die räumlichen Beobachtungen Realisierungen 
eines Zufallsfeldes sind.\\
Im Folgenden bezeichne $(\Omega, \mathcal{F}, \mathbb{P})$ einen Wahrscheinlichkeitsraum, wobei
\begin{itemize}[topsep=2pt,parsep=1pt,itemsep=0.5pt]
\item[$\circ$] $\Omega$ eine beliebige Menge 
\item[$\circ$] $\mathcal{F}$ eine $\sigma$-Algebra auf $\Omega$
\item[$\circ$] $\mathbb{P}$ ein Wahrscheinlichkeitsmaß 
auf $\mathcal{F}$
\end{itemize}
bezeichnet.

Des Weiteren sei $E$ ein vollständiger und separabler metrischer Raum
mit Borel'scher $\sigma$-Algebra $\mathcal{B}(E)$ und 
$\mathcal{S} \subset \mathbb{R}^d$ eine abzählbare Teilmenge 
mit $d \in\mathbb{N},\ d \geq 2$.

\begin{definition} \textbf{Zufallsfeld} \\
Eine Familie von Zufallsvariablen 
$\mathrm{Z}_\mathcal{S}:=\lbrace \mathrm{Z}(s),\ s \in \mathcal{S}\rbrace$
auf $(­\Omega, \mathcal{F},\mathbb{P})$ heißt
\underline{Zufallsfeld} oder \underline{stochastischer Prozess}.
Insbesondere sind die einzelnen Zufallsvariablen messbare Abbildungen:
$\mathrm{Z}: \Omega \times \mathcal{S} \longrightarrow \mathbb{R}$
(vgl.\,z.\,B.\,\cite[S.\,322]{Spodarev2015}, \cite[S.\,7]{Adler2007}).
\end{definition}  

Es ist wichtig zu betonen, dass die einzelnen Zufallsvariablen 
$\mathrm{Z}(s)$, die ein Zufallsfeld $\mathrm{Z}_\mathcal{S}$ konstituieren,
\textit{nicht stochastisch unabhängig} sind, sondern entsprechend dem
Aspekt der Verbundenheit des Tobler'schen Gesetzes (vgl.\,Punkt c) der
Auflistung in Kap.\,\ref{KapitelTheorieKernidee} oben) sehr wohl miteinander 
in Beziehung stehen.\\

Durch die folgende Einschränkung ergibt sich aus der Definition eines 
Zufallsfeldes die Definition eines Gauß'schen Zufallsfeldes:

\begin{definition} \textbf{Gauß'sches Zufallsfeld} \\
Ein Zufallsfeld  
$\mathrm{Z}_\mathcal{S}$ auf 
einem Wahrscheinlichkeitsraum $(\Omega, \mathcal{F}, \mathbb{P})$ heißt
\underline{Gauß'sches Zufallsfeld}\footnote{Gauß`sche Zufallsfelder werden 
in Lehrbüchern der Stochastik synonym auch als Gauß-Prozesse bezeichnet 
(vgl.\,z.\,B.\,\cite[S.\,462]{Schmidt2011}, 
\cite[S.\,344]{Meintrup2005}).\\}, 
wenn für \textit{alle} endlichen, 
nicht-leeren Teilmengen $\mathcal{J} \subset \mathcal{S}$ der Mächtigkeit 
$\#\mathcal{J} = n \in \mathbb{N}$ der Zufallsvektor 
\begin{equation}
\label{Gl_Def_Gaußfeld_1}
\mathrm{Z}_\mathcal{J} :=\left( \mathrm{Z}_j\right)_{j \in \mathcal{J}}=
\left(\mathrm{Z}_{j_1}, \mathrm{Z}_{j_2}, ..., \mathrm{Z}_{j_n}\right) 
\end{equation}
eine multivariate Normal-Verteilung besitzt 
(vgl.\,z.\,B.\,\cite[S.\,46f]{Diggle2007}, \cite[S.\,13]{Montero2015}).
Ausgeschrieben ist Gl.\,(\ref{Gl_Def_Gaußfeld_1}) gleichbedeutend zu
\begin{equation}
\label{Gl_Def_Gaußfeld_2}
\left(\mathrm{Z}_{j_1}, \mathrm{Z}_{j_2}, ..., \mathrm{Z}_{j_n}\right) =
\left(\mathrm{Z}(s_{j_1}), \mathrm{Z}(s_{j_2}), ..., 
\mathrm{Z}(s_{j_n})\right) \sim
\mathcal{N}_n\left(\mu, \mathbf{\Sigma}\right),
\end{equation}
mit den bekannten Parametern der n-dimensionale Normal-Verteilung
$\mathcal{N}_n$
\begin{itemize}[topsep=2pt,parsep=1pt,itemsep=0.5pt]
\item[$\circ$] dem Erwartungswert-Vektor $\mu$
\item[$\circ$] einer symmetrischen, positiv semi-definiten 
Kovarianz-Matrix $\mathbf{\Sigma}$.
\end{itemize}
\end{definition}  

Nun existiert in der Realität meist nur eine einzige Realisation 
$z(s)$ der Zufallsvariablen $\mathrm{Z}(s)$. Aus nur \textit{einer}
bekannten Merkmalsausprägung je Ort können ohne weitere
 Annahmen aber weder Erwartungswert noch Kovarianz geschätzt 
werden. Dies resultiert daraus, dass es nicht möglich ist, die
wahrscheinlichkeitstheoretischen Eigenschaften
eines Zufallsfeldes zu bestimmen, wenn nur eine einzige
Realisierung vorliegt (dies wäre vergleichbar mit einer 
Schätzung basierend auf einer Stichprobe vom Umfang eins).
Um konsistente Inferenz zu ermöglichen, sind hingegen viele
Realisierungen notwendig.\\
Darüber hinaus ist selbst von dieser einen Realisierung nur ein Teil
bekannt, nämlich die Werte an den Punkten der Stichproben-Ziehung.\\

Die Lösung für diese Schwierigkeit besteht darin, geeignete
Annahmen bezüglich der räumlichen Homogenität oder
 \textit{Stationärität} zu treffen.\\ 
Die Idee hinter dieser Annahme besteht darin,
Wiederholungen der (nicht zugänglichen) Realisierungen an
\textit{einem} Ort durch Wiederholungen im Raum zu ersetzen.
Das bedeutet, dass Werte, die an verschiedenen Orten
im untersuchten Gebiet beobachtet werden, dieselben Eigenschaften
besitzen und mathematisch als Realisierungen desselben 
Zufallsfeldes betrachtet werden können.
Die Stationaritätshypothese impliziert somit, dass die 
gemeinsamen Verteilungen des Zufallsfeldes 
translationsinvariant ist.\\

Die folgenden Ausführungen zu Stationarität und Isotropie
finden sich auch z.\,B.\,in 
\cite[S.\,52-53]{Cressie1993}, \cite[S.\,16-17]{Chiles1999},
\cite[S.\,14-16]{Montero2015}, \cite[S.\,107]{Benndorf2023},
\cite[S.\,52-57]{Schmid2025}.

\begin{definition} \textbf{Strenge Stationarität} \\
Das Zufallsfeld $\mathrm{Z}_\mathcal{S}$ heißt 
\underline{streng stationär}, wenn für alle $k \in \mathbb{N}$ 
die Familien von Zufallsvariablen
\begin{equation*}
\lbrace \mathrm{Z}(s_1), \mathrm{Z}(s_2), \ldots, \mathrm{Z}(s_k)
\rbrace
\end{equation*}
und
\begin{equation*}
\lbrace \mathrm{Z}(s_1+h), \mathrm{Z}(s_2+h), \ldots,\mathrm{Z}(s_k+h)
\rbrace
\end{equation*}
für beliebige Punkte $s_1, s_2, \ldots, s_k \in \mathcal{S}$
und jeden Translationsvektor $h \in \mathbb{R}^d$,
deren verschobene Punkte noch im Gebiet $\mathcal{S}$ liegen,
dieselbe gemeinsame Verteilungsfunktion besitzen.

Anders ausgedrückt:
Die gemeinsame Verteilungsfunktion
\begin{equation*}
F_\mathrm{Z}(z)\left( \mathrm{Z}(s_1), \mathrm{Z}(s_2), \ldots, \mathrm{Z}(s_k) \right):= \mathbb{P}\left(
\mathrm{Z}(s_1) \leq z_1, \mathrm{Z}(s_2) \leq z_2, \ldots, 
\mathrm{Z}(s_k) \leq z_k\right)
\end{equation*}
der Familie $\lbrace \mathrm{Z}(s_1), \mathrm{Z}(s_2), \ldots, 
\mathrm{Z}(s_k)\rbrace$ wird durch eine Translation um eine 
beliebigen Vektor $h$ nicht verändert.
\end{definition}

Hierbei handelt es sich um eine sehr restriktive Bedingung.
Aus diesem Grund wird diese Forderung meist gelockert 
und durch die der schwachen Stationarität ersetzt, die nur die 
ersten beiden Momente des Zufallsfeldes involviert.

\begin{definition} \textbf{Schwache Stationarität,
stationäre Kovarianz-Funktion} \\
Das Zufallsfeld $\mathrm{Z}_\mathcal{S}$ heißt
\underline{schwach stationär}, wenn es die folgenden 
Bedingungen erfüllt:

\begin{itemize}[topsep=2pt,parsep=1pt,itemsep=0.5pt]
\item[$\circ$] Der Erwartungswert existiert und ist konstant; 
er hängt also nicht vom Ort $s$ ab:
\begin{equation}
\label{Gl_Def_schwach_stationär_Erwartungswert}
\mathbb{E}\left( \mathrm{Z}(s) \right) = \mu(s) = \mu.
\end{equation}
\item[$\circ$] Die Kovarianz existiert für jedes Paar von
Zufallsvariablen $\mathrm{Z}(s)$ und $\mathrm{Z}(s+h)$ und 
hängt \textit{nur} vom Verschiebungsvektor $h$ ab, der die 
Orte $s$ und $(s+h)$ verbindet, jedoch nicht von deren 
konkreten Lage innerhalb des betrachteten Raums $\mathcal{S}$:
\begin{equation}
\label{Gl_Def_schwach_stationär_Kovarianz}
\mathrm{Cov} \left( \mathrm{Z}(s), \mathrm{Z}(s+h) \right) =:
C(h),\quad \forall s \in \mathcal{S} \textrm{ und } h \in \mathbb{R}^d.
\end{equation}
\end{itemize}
Die Funktion $C(\cdot): \mathbb{R}^d \rightarrow \mathbb{R}$ 
wird als \underline{stationäre Kovarianz-Funktion} bezeichnet.
\end{definition}

Da nach Voraussetzung die Kovarianz eines schwach stationären 
Zufallsfeldes nur vom Verschiebungsvektor $h$ abhängt, folgt daraus, 
dass die Varianz eines solchen Zufallsfeldes existiert und konstant 
ist, denn:
\begin{equation}
\label{Gl_VarianzZF_schwach_stationär}
\mathrm{Var} \left( \mathrm{Z}(s) \right) = 
\mathrm{Cov} \left( \mathrm{Z}(s), \mathrm{Z}(s) \right) = C(0) =: 
\sigma^2 \in \mathbb{R}_+\setminus \lbrace 0 \rbrace.
\end{equation}

\begin{definition} \textbf{Isotropie/Anisotropie} \\
Falls die Kovarianz-Funktion $C(\cdot)$ eines stationären Zufallsfeld 
rotationsinvariant ist, d.\,h.\,nur von der Distanz 
$\Vert s_1 - s_2\Vert$ zwischen den betrachteten 
Punkten $s_1$ und $s_2$ abhängt aber \textit{nicht} von der Richtung 
des Differenzvektors $(s_1 - s_2)$, dann wird das
Zufallsfeld als \underline{isotrop} bezeichnet, andernfalls 
als \underline{anisotrop}.
\end{definition}

Unter bestimmten Umständen lässt sich der Fall eines anisotropen
Zufallsfeldes durch Koordinaten-Transformation auf den eines 
isotropen Zufallsfeldes zurückführen. Die Handhabung von Anisotropie wird detailliert in Kap.\,\ref{KapitelTheorieAnisotropie} behandelt.\\

Kovarianzfunktionen beschreiben die räumliche Abhängigkeit vollständig.
In der Praxis arbeitet die räumlichem Statistik meist mit der eng 
verwandten Semivarianz und deren grafischen Darstellung, dem Variogramm.  